% !TeX root = ../Book1.tex
% 外法向优先约定在环域外圈和内圈上的不同方向。
\begin{figure}[htbp]
	\centering
	\begin{tikzpicture}[font=\small,scale=1.05]
		\fill[GeometryBlueLight,even odd rule] (0,0) circle[radius=2] (0,0) circle[radius=0.88];
		\draw[GeometryBlue,thick] (0,0) circle[radius=2];
		\draw[GeometryBlue,thick] (0,0) circle[radius=0.88];
		\draw[GeometryBlue,very thick,->] (110:2) arc[start angle=110,end angle=155,radius=2];
		\draw[GeometryBlue,very thick,->] (245:2) arc[start angle=245,end angle=290,radius=2];
		\draw[GeometryBlue,very thick,->] (150:0.88) arc[start angle=150,end angle=95,radius=0.88];
		\draw[GeometryBlue,very thick,->] (330:0.88) arc[start angle=330,end angle=275,radius=0.88];
		\fill[GeometryBlue] (2,0) circle[radius=1.6pt];
		\draw[orange!75!black,thick,->] (2,0)--(2.95,0) node[right] {\(\nu\)};
		\draw[GeometryBlue,thick,->] (2,0)--(2,0.85) node[above] {\(T\)};
		\fill[GeometryBlue] (0,-0.88) circle[radius=1.6pt];
		\draw[orange!75!black,thick,->] (0,-0.88)--(0,-0.15) node[right] {\(\nu\)};
		\draw[GeometryBlue,thick,->] (0,-0.88)--(-0.72,-0.88) node[below] {\(T\)};
		\node[GeometryGray] at (-0.05,0.4) {圆孔};
		\node[GeometryBlue] at (0,1.45) {环域 \(A\)};
		\node[GeometryGray] at (-2.8,1.4) {外圈逆时针};
		\node[GeometryGray] at (-2.8,-1.4) {内圈顺时针};
	\end{tikzpicture}
	\caption{平面取标准定向时，外法向优先要求 \((\nu,T)\) 为正向。沿边界正向 \(T\) 行进，环域总在左侧；内圈的外侧是圆孔，因此内圈与外圈的绕行方向相反。图中的外向量不要求单位长度。}
	\label{g1:fig:boundary-outward-first}
\end{figure}
